
\label{subsec:metadata}

A key goal in expanding the context provided to the model is to train on a broader, more diverse dataset of trajectories. Instead of just using high-quality demonstration data, \MethodName{} leverages lower quality demonstrations (including failures) and even autonomous data from prior models. Since we still want \MethodName{} to perform the task as well as possible at test time, we need to appropriately label these diverse trajectories with information about \emph{how} the task was performed so that the model can correctly contextualize them. To this end, we add a variety of ``episode metadata'' information to the context with attributes of the given training episode. We denote the set of metadata $m$, which may contain various labels including



\begin{itemize}[leftmargin=10pt]
    \item \textbf{Overall speed}: the length of the episode in timesteps. We discretize the values in an interval of 500 steps, i.e., values between 1750 and 2250 are binned to ``2000 steps''. Often faster speed also corresponds to higher quality, e.g., the episode has fewer mistakes.
    \item \textbf{Overall quality}: task execution quality expressed as a score between 1 and 5, with 5 being the highest quality.
    \item \textbf{Mistake}: label indicating whether the robot made a mistake within a given action segment (e.g., failing to grasp an object or performing the wrong subtask). These labels are provided by humans coarsely annotating our data.
\end{itemize}

The \MethodName{} model is thus trained with ground-truth episode speed and manual annotations of the episode quality and mistake segments from a diverse data mixture. The data diversity (e.g., episodes of varying speed) provides the necessary signals for the model to learn to correlate such metadata with the target action.
At runtime, the model can then be instructed to perform the task at high speed, with high quality, and without mistakes, through metadata prompting.

